% ================================================================ TD 2
\td{2}{Cluster Docker Swarm}

J'ai monté un cluster Docker Swarm local de 4 machines (1 manager et 3 nœuds), chacune étant un conteneur Docker-in-Docker.

% ---------------------------------------------------------------- prérequis
\prerequis{images publiées sur un registry}

\subsection{Publier les images avec GitHub Actions}

Swarm ne sait pas construire d'images : j'ai donc écrit deux workflows qui les construisent à chaque push sur \code{main} et les publient sur GHCR.

\begin{fichier}{.github/workflows/build.server.yml}
env:
  REGISTRY: ghcr.io
  IMAGE_NAME: ${{ github.repository }}/server
\end{fichier}

\subsection{Vérifier les images sur GHCR}

\capture[0.20\textheight]{td2-00-ghcr-packages.png}{Les deux images sont publiques : les nœuds les téléchargent sans connexion.}

% ---------------------------------------------------------------- exercice 2
\exercice{2}{Création du cluster Docker Swarm}

\subsection{Lancer les 4 machines}

Les nœuds \code{node1} à \code{node3} sont définis comme le manager, et \code{privileged: true} permet de faire tourner Docker dans un conteneur.

\begin{fichier}{swarm/compose.yml}
services:
  manager:
    image: docker:dind
    hostname: manager
    privileged: true
  # ...
\end{fichier}

\begin{commandes}
cd swarm
docker compose up -d
docker ps
\end{commandes}

\capture[0.34\textheight]{td2-01-dind-up.png}{Les 4 machines du cluster tournent à côté des conteneurs du TD 1.}

\begin{commandes}
docker compose exec manager docker ps
docker compose exec node1 docker ps
\end{commandes}

La liste est vide dans les deux cas : chaque machine a son propre Docker, séparé de celui du PC.

\subsection{Initialiser le cluster}

\begin{commandes}
docker compose exec manager docker swarm init
\end{commandes}

\subsection{Faire rejoindre les nœuds}

\begin{commandes}
TOKEN=$(docker compose exec manager docker swarm join-token -q worker)
docker compose exec node1 docker swarm join --token $TOKEN manager:2377
docker compose exec node2 docker swarm join --token $TOKEN manager:2377
docker compose exec node3 docker swarm join --token $TOKEN manager:2377
\end{commandes}

\subsection{Lister les nœuds du cluster}

\begin{commandes}
docker compose exec manager ash
docker node ls
exit
\end{commandes}

\capture[0.25\textheight]{td2-05-node-ls.png}{Les 4 nœuds sont \code{Ready}, avec le \code{manager} comme \code{Leader}.}

% ---------------------------------------------------------------- exercice 3
\exercice{3}{Tests du cluster}

\subsection{Créer une stack hello-world}

\begin{fichier}{swarm/hello-world.compose.yml}
    # ...
    deploy:
      replicas: 2
\end{fichier}

\subsection{Copier la stack dans le manager}

J'ai ajouté un volume anonyme sur \code{/home/manager} pour que la stack copiée survive à une recréation du conteneur.

\begin{fichier}{swarm/compose.yml}
  manager:
    # ...
    volumes:
      - /home/manager
\end{fichier}

\begin{commandes}
docker compose up -d manager
docker compose exec manager docker node ls
docker compose exec manager mkdir -p /home/manager
docker compose cp hello-world.compose.yml manager:/home/manager/
docker compose exec manager ls /home/manager
\end{commandes}

Après sa recréation, le manager voit toujours ses 4 nœuds \code{Ready}.

\subsection{Déployer la stack et observer le cluster}

\subsubsection{Relancer le cluster après un redémarrage du PC}

\begin{commandes}
docker compose start
docker compose exec manager docker node ls
\end{commandes}

\probleme{Swarm retient l'IP du manager, qui avait changé : les nœuds restaient \code{Down}.}{J'ai fait quitter le cluster à toutes les machines, puis je l'ai recréé.}

\begin{commandes}
docker compose exec node1 docker swarm leave --force
docker compose exec node2 docker swarm leave --force
docker compose exec node3 docker swarm leave --force
docker compose exec manager docker swarm leave --force

docker compose exec manager docker swarm init
TOKEN=$(docker compose exec manager docker swarm join-token -q worker)
docker compose exec node1 docker swarm join --token $TOKEN manager:2377
docker compose exec node2 docker swarm join --token $TOKEN manager:2377
docker compose exec node3 docker swarm join --token $TOKEN manager:2377

docker compose exec manager docker node ls
\end{commandes}

\capture{td2-08-swarm-rebuild.png}{Le cluster recréé retrouve ses 4 nœuds \code{Ready}.}

\subsubsection{Déployer la stack hello}

\begin{commandes}
docker compose exec manager docker stack deploy -c /home/manager/hello-world.compose.yml hello
docker compose exec manager docker service ls
docker compose exec manager docker service ps hello_hello
\end{commandes}

\capture{td2-09-stack-deploy.png}{\code{hello-world} s'arrête dès son message affiché : le service reste à \code{0/2} et Swarm relance sans cesse les copies, sur tous les nœuds.}

\needspace{5cm}
Pour comparer, un service nginx qui reste allumé affiche bien \code{3/3}.

\begin{commandes}
docker compose exec manager docker service create --name web --replicas 3 nginx:alpine
docker compose exec manager docker service ls
docker compose exec manager docker service rm web

docker compose exec manager docker ps -a
docker compose exec node1 docker ps -a
docker compose exec node2 docker ps -a
docker compose exec node3 docker ps -a
\end{commandes}

\capture{td2-10-docker-ps-nodes.png}{Swarm a réparti les copies sur les 4 machines, \code{manager} compris. Chaque machine ne voit que les siennes.}

\subsection{Faire varier la clause deploy}

\subsubsection{Sur le manager uniquement}

\begin{fichier}{swarm/hello-world-manager.compose.yml}
    # ...
    deploy:
      replicas: 2
      placement:
        constraints:
          - node.role == manager
\end{fichier}

\begin{commandes}
docker compose exec manager docker stack rm hello
docker compose cp hello-world-manager.compose.yml manager:/home/manager/
docker compose exec manager docker stack deploy -c /home/manager/hello-world-manager.compose.yml hello
docker compose exec manager docker service ps hello_hello
\end{commandes}

\capture{td2-11-stack-manager.png}{Toutes les copies tournent sur le \code{manager}, aucune sur les workers.}

\subsubsection{Sur tous les nœuds (mode global)}

\begin{fichier}{swarm/hello-world-global.compose.yml}
    # ...
    deploy:
      mode: global
\end{fichier}

\begin{commandes}
docker compose exec manager docker stack rm hello
docker compose cp hello-world-global.compose.yml manager:/home/manager/
docker compose exec manager docker stack deploy -c /home/manager/hello-world-global.compose.yml hello
docker compose exec manager docker service ls
docker compose exec manager docker service ps hello_hello
\end{commandes}

\capture[0.45\textheight]{td2-12-stack-global.png}{En mode \code{global}, Swarm lance une copie sur chacun des 4 nœuds.}
